\section{Conclusion}
\label{sec:conclusion}

\textbf{Advantages.}
PatchCore needs no training: a forward pass and coreset selection fit a category from about 200 normal images in \pcFitAll{}\,s.
It is fast and light (\pcInferAll{}\,ms per image on the GPU, \pcInferEtoEAll{}\,ms including loading, and a 57--88\,MiB bank) and accurate (image AUROC \pcImgAurocAll{}, AU-PRO \pcAuproAll{}).
It is robust to pose (screw image AUROC \pcImgAurocScrew{} vs \pdImgAurocScrew{}), and every anomalous patch can be explained by showing its nearest normal patch.
PaDiM, a much simpler statistical model, performs equally well on stationary textures but requires aligned parts, about 3.8\,GB of covariances per category, and a uniform background free of specks.

\textbf{Limitations.}
The largest practical weakness is threshold calibration with few normal images.
With 21--32 validation images, the 99\% quantile is close to the validation maximum.
This collapses PaDiM's F1 on tile and screw and costs PatchCore 6.5 F1 points relative to the oracle reference.
Normal drift between training and test images causes most of PatchCore's texture false alarms.
Defects smaller than one feature cell are detected only about half the time, and thin defects are segmented as wide blobs, even though AU-PRO looks high.
The centre crop discards the image border, and every product needs its own memory bank and threshold.
Our evidence is limited to a single seed, six of the 15 categories and a 224-pixel input.

\textbf{Industrial applicability.}
PatchCore suits a line that can supply a few hundred good images per product and tolerate per-product setup, which takes seconds; inference takes milliseconds on one GPU.
Before deployment we would
(i)~collect a larger normal calibration set and choose the operating point from the costs of a miss and a false alarm;
(ii)~inspect the full field of view without cropping;
(iii)~use higher resolution or tiling for defects smaller than 0.3\% of the image;
(iv)~monitor normal score levels and refresh the bank when the material lot or lighting changes; and
(v)~send near-threshold images (within about $\pm$10\%) to human review.
Promising next steps are a native higher-resolution protocol fixed in advance, layer-1 maps for segmentation combined with the layer-2+3 score for detection, an a-priori re-test of DINOv2 without the final LayerNorm, rotation augmentation of the bank for pose-variable parts, and multiple seeds with confidence intervals.
